\chapter{Discussion} % Chapter 5
\label{Discussion}

% =============================================================================
% DISCUSSION (5-8 pages)
% Rubric: Clarity of Conclusions (10% of Report)
% =============================================================================

% -----------------------------------------------------------------------------
% SECTION 5.1: Interpretation of Results
% -----------------------------------------------------------------------------
\section{Interpretation of Results}
% TODO: The non-monotonic pattern
%   - L0 most cooperative, L1 most aggressive, L2 defensive, L3 exploitative (pilot)
%   - This is NOT "more reasoning = more cooperation" (naive expectation)
%   - This is NOT "more reasoning = more conflict" (Kuusela & Roy's finding)
%   - It's non-monotonic and mechanism-dependent
%
% TODO: Connect to K-Level Reasoning literature
%   \cite{zhang2024klevel}: diminishing returns at high depth — consistent
%   \cite{jia2025cot}: CoT interacts with reasoning level — consistent
%   \cite{duan2024gtbench}: CoT sometimes hurts — consistent
%   TMGBench: second-order ToM minimal gain — consistent with L2→L3

% TODO: The arm_cost interaction effect
%   - L0 insensitive to game structure (always cooperates ~70%)
%   - L2 extremely sensitive (Gini 0.58 → 0.84 with arm_cost change)
%   - Interpretation: deeper reasoning AMPLIFIES structural incentives
%   - L0 agents don't notice the game changed; L2 agents optimize for it
%
% TODO: Compare L2 fear spiral to Kuusela & Roy
%   - Their L2 fear spiral manifests as PRE-EMPTIVE ATTACKS (only option in 2-action game)
%   - Our L2 fear spiral manifests as DEFENSIVE HOARDING (arm_self 37%, attack only 2%)
%   - Same reasoning depth, same fear mechanism, different behavioral expression
%   - Difference is action space: our game has arm_self as intermediate option
%   - Hypothesis: if arm_self were removed, L2 would also show pre-emptive attacks
%   - Their game is also one-shot destruction (no repeated game, no recovery) —
%     agents can't move toward Pareto because a single attack = game over
%   - Our game allows recovery → repeated interaction → different equilibria possible
%   - Key point: the fear spiral from deeper reasoning is CONSERVED across game designs,
%     but the behavioral expression depends on available actions and game structure

% TODO: Connect to computational depth interpretation
%   \cite{pfau2024dots}: the manipulation works because it changes computation,
%     not just content. Even if traces are unfaithful, the prompt changes what
%     the model computes before acting.

% -----------------------------------------------------------------------------
% SECTION 5.2: The Origins of Order
% -----------------------------------------------------------------------------
\section{The Origins of Order}
% TODO: Answer the title question
%   - Order emerges from the INTERACTION of reasoning and structure
%   - Neither alone determines outcomes:
%     * Same structure, different reasoning → different outcomes (reasoning depth effect)
%     * Same reasoning, different structure → different outcomes (parameter sweeps)
%     * The interaction is what matters (factorial results)
%
% TODO: Phase transition shift (if confirmed by origins factorial)
%   - The critical radius depends on reasoning depth
%   - Deeper reasoning may LOWER the cooperation threshold
%     (or raise it — depends on results)
%   - This is the strongest theoretical contribution

% -----------------------------------------------------------------------------
% SECTION 5.3: Faithfulness and Causality
% -----------------------------------------------------------------------------
\section{On Reasoning Traces}
% TODO: What can we claim about traces?
%   - Traces correlate with behavior (L2 mentions opponent modeling, L0 doesn't)
%   - But correlation ≠ causation in trace content
%   - The PROMPT manipulation is causal (it's the experimental variable)
%   - The trace CONTENT may or may not reflect the actual computation
%
% TODO: Connect to Thought Anchors/Branches
%   \cite{bogdan2025anchors}: sentence-level causal impact possible via resampling
%   \cite{bogdan2025branches}: need distributions, not single traces
%   - Limitation: we don't do full resampling (compute prohibitive at 30 agents)
%   - Future work: apply Thought Anchors to critical rounds

% TODO: The "computation, not content" framing
%   - Our L0→L3 prompts change computational depth
%   - This is a real intervention even if traces are unfaithful
%   - Frame results as: "reasoning depth changes behavior" not "reasoning
%     content explains behavior"
%
% TODO: Design choice — instruction, not computation budget
%   - We manipulate instructed reasoning depth, NOT computational budget
%   - Token counts (L0=6, L1=31, L2=40, L3=44) are OBSERVED, not controlled
%   - L2 and L3 produce similar computation (~40-44 tokens) but qualitatively
%     different behavior (37% arm vs 8% attack) — the instruction matters
%     beyond just token count
%   - We deliberately do NOT force longer reasoning for higher levels:
%     that would confound instruction with computation length
%   - Tension: \cite{pfau2024dots} argues more tokens = more compute = better.
%     \cite{lanham2023faithfulness} argues more tokens = more rationalization.
%     Our design does not resolve this — it's a single manipulation (instruction)
%     with token count as observed mediator, not as intervention.
%   - This is a strength: clean causal claim (prompt → action), not overclaim

% -----------------------------------------------------------------------------
% SECTION 5.4: Limitations
% -----------------------------------------------------------------------------
\section{Limitations}
% TODO: Model specificity
%   - Results on Gemma 2 27B — may not generalize to other architectures
%   - Robustness check on second model partially addresses this
%
% TODO: Game specificity
%   - One coordination game — different games may show different patterns
%   - But game structure was systematically varied
%
% TODO: Prompt sensitivity
%   - Exact wording of reasoning instructions matters
%   - We chose principled levels mapped to K-Level framework
%   - But other phrasings might produce different results
%   - Report FormatSpread (\cite{sclar2024}) or PromptSensiScore (\cite{zhuo2024})
%   \cite{lore2024}: showed framing effects dominate in LLM games — we must
%     distinguish reasoning depth effects from mere framing effects
%
% TODO: Scale
%   - 30 agents — large enough for structure but not "society-scale"
%   - Computational constraints limit scale
%
% TODO: Trace faithfulness
%   - Cannot fully verify traces reflect internal computation
%   - We frame claims carefully: computational depth, not cognitive claims
%
% TODO: Zero-cost regime discussion
%   - Initial experiments used zero arm/conflict costs
%   - This makes some actions trivially dominant
%   - Cost sweeps address this, but base param choice affects results

% -----------------------------------------------------------------------------
% SECTION 5.5: Implications
% -----------------------------------------------------------------------------
\section{Implications}
% TODO: For computational social science
%   - LLM agents are NOT neutral simulation tools — their reasoning shapes outcomes
%   - Researchers using LLM-ABMs must control for reasoning depth
%   \cite{larooij2025}: validation is central challenge
%   \cite{taillandier2025}: hybrid approaches needed for transparency
%   \cite{riedl2025}: prompt design steers emergent coordination — consistent
%
% TODO: For AI safety/alignment
%   - Deeper reasoning doesn't mean "better" social outcomes
%   - Non-monotonic effects: L2 agents arm more, L3 agents exploit more
%   - "More capable" AI may produce MORE inequality, not less
%
% TODO: For game theory
%   - Level-K reasoning in LLMs behaves differently than in RL
%   - Interaction with environment structure is key
%   \cite{lore2024}: framing > game structure — our finding extends this by
%     showing reasoning depth ALSO shapes outcomes, and interacts with structure
%   \cite{pellert2025}: LLMs can replicate human cooperation patterns, but our
%     finding shows this depends on reasoning depth instruction
